\chapter{Moduli space of curves}
\label{chap:moduli}
% archon:covers MR4276287UniformityInMordellLangForCurves/ModuliSpace.lean

This chapter establishes the fine moduli space \(\mathbb{M}_g\) of smooth genus-\(g\) curves
with level structure, together with the universal curve over it and the Torelli map to the
moduli of principally polarized abelian varieties. These objects provide the ambient setting
for \cref{def:ambient_setup}.

\section{Level structure and the fine moduli space}

\begin{definition}[Level structure]
  \label{def:level_structure}
  \lean{MR4276287UniformityInMordellLangForCurves.ModuliSpace.levelStructure}
  % SOURCE: Section 6 of the paper, subsec:univfamily (read from references/MR4276287-mordell-lang-curves-source/SettingUp.tex)
  % SOURCE QUOTE: "$\mathbb{M}_g$ denotes the fine moduli space of smooth curves of genus $g$, with level-$\ell$-structure where $\ell\ge \bestlevel$ is fixed, \textit{cf.} \cite[Chapter~XVI, Theorem~2.11 (or above Proposition~2.8)]{ACG:Curve}, \cite[(5.14)]{DM:irreducibility}, or \cite[Theorem 1.8]{OortSteenbrink}."
  \textit{Source: Section 6 of \cite{MR4276287-mordell-lang-curves} (read from references/MR4276287-mordell-lang-curves-source/SettingUp.tex); formal definition in \cite[Theorem~1.8]{OortSteenbrink}.}
  Fix integers \(g \ge 2\) and \(\ell \ge 3\). A \emph{level-\(\ell\) structure} on a smooth
  projective curve \(C\) of genus \(g\) is an isomorphism of group schemes
  \[
    \mathrm{Jac}(C)[\ell] \xrightarrow{\;\sim\;} (\mathbb{Z}/\ell\mathbb{Z})^{2g}
  \]
  compatible with the Weil pairing on \(\mathrm{Jac}(C)[\ell]\). Here \(\mathrm{Jac}(C)[\ell]\)
  denotes the \(\ell\)-torsion subgroup scheme of the Jacobian.
\end{definition}

\begin{definition}[Fine moduli space with level structure]
  \label{def:Mg_level}
  \lean{MR4276287UniformityInMordellLangForCurves.ModuliSpace.mgLevel}
  \uses{def:level_structure}
  % SOURCE: Section 6 of the paper, subsec:univfamily (read from references/MR4276287-mordell-lang-curves-source/SettingUp.tex)
  % SOURCE QUOTE: "Recall from $\mathsection$\ref{subsec:hgtineqnondeg} that $\mathbb{M}_g$ denotes the fine moduli space of smooth curves of genus $g$, with level-$\ell$-structure where $\ell\ge \bestlevel$ is fixed, \textit{cf.} \cite[Chapter~XVI, Theorem~2.11 (or above Proposition~2.8)]{ACG:Curve}, \cite[(5.14)]{DM:irreducibility}, or \cite[Theorem 1.8]{OortSteenbrink}."
  \textit{Source: Section 6 of \cite{MR4276287-mordell-lang-curves} (read from references/MR4276287-mordell-lang-curves-source/SettingUp.tex).}
  Fix \(g \ge 2\) and \(\ell \ge 3\). The \emph{fine moduli space} \(\mathbb{M}_g\) of smooth
  genus-\(g\) curves with level-\(\ell\) structure is a scheme over
  \(\mathrm{Spec}\,\mathbb{Z}[1/\ell]\) that represents the moduli functor
  associating to each \(\mathrm{Spec}\,\mathbb{Z}[1/\ell]\)-scheme \(T\) the set of
  isomorphism classes of pairs \((C \to T, \alpha)\), where \(C \to T\) is a smooth proper
  curve of genus \(g\) and \(\alpha\) is a level-\(\ell\) structure on \(C\). For \(\ell \ge 3\),
  this functor is representable, so \(\mathbb{M}_g\) is a fine (not merely coarse) moduli scheme.
\end{definition}

\begin{definition}[Universal curve]
  \label{def:universal_curve}
  \lean{MR4276287UniformityInMordellLangForCurves.ModuliSpace.universalCurve}
  \uses{def:Mg_level}
  % SOURCE: Section 6 of the paper, subsec:univfamily (read from references/MR4276287-mordell-lang-curves-source/SettingUp.tex)
  % SOURCE QUOTE: "There exists a universal curve $\mathfrak{C}_g$ over $\mathbb{M}_g$, it is smooth and proper over $\mathbb{M}_g$ and its fibers are smooth curves of genus $g$. Moreover, $\mathfrak{C}_g\rightarrow \mathbb{M}_g$ is projective, \textit{cf.} \cite[Corollary to Theorem 1.2]{DM:irreducibility} or \cite[Remark 2, \S 9.3]{NeronModels}."
  \textit{Source: Section 6 of \cite{MR4276287-mordell-lang-curves} (read from references/MR4276287-mordell-lang-curves-source/SettingUp.tex).}
  There exists a \emph{universal curve} \(\mathfrak{C}_g \to \mathbb{M}_g\): a smooth proper
  projective morphism of relative dimension \(1\) whose geometric fibers are smooth curves of
  genus \(g\). By the fine moduli property of \(\mathbb{M}_g\), the pair
  \((\mathbb{M}_g, \mathfrak{C}_g)\) is initial among all families of smooth genus-\(g\) curves
  with level-\(\ell\) structure.
\end{definition}

\section{Geometric properties}

\begin{proposition}[Quasi-projectivity of \(\mathbb{M}_g\)]
  \label{prop:Mg_quasiprojective}
  \lean{MR4276287UniformityInMordellLangForCurves.ModuliSpace.mgQuasiProjective}
  \uses{def:Mg_level}
  % SOURCE: Section 6 of the paper, subsec:univfamily (read from references/MR4276287-mordell-lang-curves-source/SettingUp.tex)
  % SOURCE QUOTE: "It is known that $\mathbb{M}_g$ is a regular, quasi-projective variety of dimension $3g-3$. We regard it over $\IQbar$; it is irreducible according to our convention introduced below Theorem~\ref{ThmHtInequality}."
  \textit{Source: Section 6 of \cite{MR4276287-mordell-lang-curves} (read from references/MR4276287-mordell-lang-curves-source/SettingUp.tex).}
  For \(\ell \ge 3\), \(\mathbb{M}_g\) is a smooth, irreducible, quasi-projective variety over
  \(\overline{\mathbb{Q}}\) of dimension \(3g - 3\). It admits a compactification
  \(\overline{\mathbb{M}_g}\) (the Deligne--Mumford compactification by stable curves with
  level structure), which is a projective variety containing \(\mathbb{M}_g\) as a Zariski
  open dense subset.
\end{proposition}

% SOURCE QUOTE PROOF: "It is known that $\mathbb{M}_g$ is a regular, quasi-projective variety of dimension $3g-3$. We regard it over $\IQbar$; it is irreducible according to our convention introduced below Theorem~\ref{ThmHtInequality}." (read from references/MR4276287-mordell-lang-curves-source/SettingUp.tex)
\begin{proof}
  \uses{def:Mg_level}
  The quasi-projectivity and regularity of \(\mathbb{M}_g\) are classical; see
  \cite[Chapter~XVI, Theorem~2.11]{ACG:Curve}, \cite[(5.14)]{DM:irreducibility}, or
  \cite[Theorem~1.8]{OortSteenbrink}. The dimension count \(\dim \mathbb{M}_g = 3g - 3\)
  is the Riemann count: a smooth curve of genus \(g \ge 2\) has \(3g - 3\) moduli.
  Irreducibility of \(\mathbb{M}_g\) over \(\overline{\mathbb{Q}}\) is a consequence of
  \cite[Theorem~5.15]{DM:irreducibility}. The Deligne--Mumford compactification
  \(\overline{\mathbb{M}_g}\) is constructed in \cite[Theorem~5.2]{DM:irreducibility} by
  allowing level structure on stable curves; it is projective by the valuative criterion and
  the properness of the moduli functor for stable curves.
\end{proof}

\section{The Torelli map}

\begin{lemma}[Torelli map is quasi-finite]
  \label{lem:torelli_quasi_finite}
  \lean{MR4276287UniformityInMordellLangForCurves.ModuliSpace.torelliQuasiFinite}
  \uses{def:Mg_level, def:universal_curve}
  % SOURCE: Section 6 of the paper, subsec:univfamily (read from references/MR4276287-mordell-lang-curves-source/SettingUp.tex)
  % SOURCE QUOTE: "As we have level structure, the Torelli morphism need not be injective on $\IC$-points, but it is finite-to-one on such points, \textit{cf.} \cite[Lemma 1.11]{OortSteenbrink}."
  \textit{Source: Section 6 of \cite{MR4276287-mordell-lang-curves} (read from references/MR4276287-mordell-lang-curves-source/SettingUp.tex).}
  The Torelli map \(\tau \colon \mathbb{M}_g \to \mathbb{A}_g\), sending a genus-\(g\) curve
  \([C]\) to its principally polarized Jacobian \([\mathrm{Jac}(C)]\), is quasi-finite. Here
  \(\mathbb{A}_g\) denotes the fine moduli space of principally polarized abelian varieties of
  dimension \(g\) with level-\(\ell\) structure.
\end{lemma}

% SOURCE QUOTE PROOF: "The Torelli map $\tau$ is quasi-finite; see \cite[Lemma 1.11]{OortSteenbrink}." (read from references/MR4276287-mordell-lang-curves-source/SettingUp.tex)
\begin{proof}
  \uses{def:Mg_level, def:universal_curve}
  By \cite[Lemma~1.11]{OortSteenbrink}, the Torelli map \(\tau\) is quasi-finite. More
  precisely, with level-\(\ell\) structure for \(\ell \ge 3\), the classical Torelli theorem
  (which asserts injectivity on \(\mathbb{C}\)-points for the coarse moduli space) implies
  that the fibers of \(\tau\) over \(\mathbb{C}\)-points are finite. Combined with the fact
  that \(\tau\) is of finite type (as a morphism between quasi-projective varieties of finite
  type over \(\overline{\mathbb{Q}}\)), quasi-finiteness follows.
\end{proof}
